\begin{table}[htbp]
\centering\small
\caption{Checks of the transient-penalty bounds on random piecewise-linear envelopes. The non-increasing envelopes have seven levels on random knots, lags drawn uniformly from $[0.02,0.4]$\,s, and target speeds from $[5,40]$\,m/s. The general envelopes have eight random levels and one lag drawn from $[0.05,0.4]$\,s. The last two columns give the median and the largest ratio of the penalty to the stated bound. The last row compares the general envelopes with the bound $\tau$ that holds for non-increasing envelopes.}\label{tab:app-transient}
\begin{tabular}{lccccc}
\toprule
Envelope & Lags & Cases & Bound & Median ratio & Largest ratio\\
\midrule
Non-increasing & $1$ & $18$ & $\sum_j\tau_j$ & $0.881$ & $0.987$\\
Non-increasing & $2$ & $17$ & $\sum_j\tau_j$ & $0.832$ & $0.986$\\
Non-increasing & $3$ & $25$ & $\sum_j\tau_j$ & $0.869$ & $0.971$\\
General & $1$ & $40$ & $\tau(1+\operatorname{Var}^+)$ & $0.644$ & $0.903$\\
General & $1$ & $40$ & $\tau$ & $2.314$ & $4.033$\\
\bottomrule
\end{tabular}
\end{table}

\begin{table}[htbp]
\centering\small
\caption{Penalty of a single lag $\tau=0.2$\,s on a doubling staircase envelope, in units of $\tau$, with each level held for about $5\tau$. The prediction is the limit of the staircase analysis. The bound is the value $\tau(1+\operatorname{Var}^+)$ of the variation bound, which is proved for Lipschitz envelopes and is shown here for comparison.}\label{tab:app-staircase}
\begin{tabular}{cccc}
\toprule
Levels & Penalty & Prediction & Bound\\
\midrule
$1$ & $1.00$ & $1.00$ & $1.00$\\
$2$ & $1.50$ & $1.50$ & $1.69$\\
$4$ & $2.49$ & $2.50$ & $3.08$\\
$6$ & $3.49$ & $3.50$ & $4.47$\\
$8$ & $4.49$ & $4.50$ & $5.85$\\
\bottomrule
\end{tabular}
\end{table}

\begin{table}[htbp]
\centering\small
\caption{The $0$--$60$\,mph penalty of the three vehicles under constant lags (s), with the time $T_{\mathrm{phys}}$ along the envelope and the penalty as a share of $\sum_j\tau_j$. The positive variation of $\ln\amax$ from rest to $60$\,mph, computed on $4001$ speeds, is $0$ for every vehicle. These runs follow the envelope without the interruptions of gear shifts, which is why the time of the Civic is shorter than in the zero-to-sixty table.}\label{tab:app-transient-veh}
\begin{tabular}{lcccc}
\toprule
Vehicle & Lags (s) & $T_{\mathrm{phys}}$ (s) & Penalty (s) & Share\\
\midrule
Leaf & $0.10$ & $10.665$ & $0.100$ & $99.6\%$\\
Leaf & $0.15, 0.05$ & $10.665$ & $0.199$ & $99.6\%$\\
Leaf & $0.30, 0.10, 0.05$ & $10.665$ & $0.448$ & $99.5\%$\\
Civic & $0.10$ & $8.949$ & $0.100$ & $99.6\%$\\
Civic & $0.15, 0.05$ & $8.949$ & $0.199$ & $99.6\%$\\
Civic & $0.30, 0.10, 0.05$ & $8.949$ & $0.448$ & $99.5\%$\\
ZR1 & $0.10$ & $3.804$ & $0.100$ & $99.6\%$\\
ZR1 & $0.15, 0.05$ & $3.804$ & $0.199$ & $99.4\%$\\
ZR1 & $0.30, 0.10, 0.05$ & $3.804$ & $0.446$ & $99.1\%$\\
\bottomrule
\end{tabular}
\end{table}

\begin{table}[htbp]
\centering\small
\caption{The $0$--$60$\,mph penalty (s) with tyre relaxation length $0.5$\,m and a powertrain lag of $0.10$\,s, against the floor speed $v_0$ of the relaxation model.}\label{tab:app-relax}
\begin{tabular}{lcccc}
\toprule
Vehicle & $v_0=2$\,m/s & $v_0=1$\,m/s & $v_0=0.5$\,m/s & $v_0=0.25$\,m/s\\
\midrule
Leaf & $0.347$ & $0.542$ & $0.778$ & $1.030$\\
Civic & $0.346$ & $0.541$ & $0.778$ & $1.029$\\
ZR1 & $0.337$ & $0.491$ & $0.664$ & $0.843$\\
\bottomrule
\end{tabular}
\end{table}
